% Solution: Tadelis (2013) Exercise 4.1 "Prove Proposition 4.1". Input by \solitem.
% New in the 50-item list. Check: check_A1.py.
\begin{exsol}{Tadelis Exercise 4.1}{76}{Exam}{2}
\asked Prove that a strictly dominant strategy equilibrium, when it exists, is the only dominant
strategy equilibrium.

\method Definitions only: strict dominance (Tadelis Def.~4.1, p.~60), strictly dominant strategy
(Def.~4.2, p.~61), strict dominant strategy equilibrium (Def.~4.3, p.~61) and Proposition~4.1
itself (p.~61); in the slides, Lecture~1, slides~23 and~25. The proof is by contradiction, one
player at a time. The statement says ``the unique dominant strategy equilibrium'' without the word
``strict'', so the proof also rules out a second profile of \emph{weakly} dominant strategies
(Lecture~1, slide~23), which is the stronger reading.

\solution
\textbf{Notation.} $\Gamma = \langle N, \{S_i\}_{i=1}^n, \{v_i\}_{i=1}^n\rangle$; for a profile
$s$, $s_{-i}$ is the profile of the others. Every $S_j$ is non-empty, so $S_{-i}$ is non-empty.
$s_i$ is \emph{strictly dominant} if $v_i(s_i, s_{-i}) > v_i(s'_i, s_{-i})$ for all
$s'_i \ne s_i$ and all $s_{-i}$; it is \emph{weakly dominant} if $v_i(s_i, s_{-i}) \ge
v_i(s'_i, s_{-i})$ for all $s'_i$ and all $s_{-i}$, with strict inequality somewhere for each
$s'_i \ne s_i$. A (strict, weak) dominant strategy equilibrium is a profile in which every player
uses a (strictly, weakly) dominant strategy.

\textbf{Step 1: a player has at most one dominant strategy when one of them is strictly
dominant.} Let $s_i^D$ be strictly dominant for player~$i$ and let $t_i \ne s_i^D$. By
Def.~4.2, for every $s_{-i} \in S_{-i}$,
\[
  v_i(s_i^D, s_{-i}) > v_i(t_i, s_{-i}). \tag{1}
\]
Suppose $t_i$ were also dominant, strictly or weakly. Then, in particular,
\[
  v_i(t_i, s_{-i}) \ge v_i(s_i^D, s_{-i}) \quad \text{for every } s_{-i} \in S_{-i}. \tag{2}
\]
Pick any $\hat s_{-i} \in S_{-i}$ (possible because $S_{-i} \ne \emptyset$). Then (1) and (2) give
$v_i(s_i^D, \hat s_{-i}) > v_i(t_i, \hat s_{-i}) \ge v_i(s_i^D, \hat s_{-i})$, a contradiction.
So no $t_i \ne s_i^D$ is dominant in either sense.

\textbf{Step 2: the proposition.} Let $s^D$ be a strict dominant strategy equilibrium, and suppose
some profile $s' \ne s^D$ is also a dominant strategy equilibrium. Since $s' \ne s^D$, there is a
player $i$ with $s'_i \ne s_i^D$. In $s'$, player~$i$'s strategy $s'_i$ is dominant; but $s_i^D$ is
strictly dominant for $i$ (because $s^D$ is a strict dominant strategy equilibrium), so by Step~1
no strategy other than $s_i^D$ can be dominant for $i$. This contradicts $s'_i \ne s_i^D$. Hence
$s^D$ is the unique dominant strategy equilibrium. $\blacksquare$

\textbf{Two remarks that complete the picture.}
\begin{itemize}
  \item The proposition does not claim existence: Battle of the Sexes has no dominant strategy
  equilibrium (Tadelis p.~62).
  \item Uniqueness extends to the other solution concepts of the course: $s^D$ is the unique
  profile that survives IESDS (Lecture~1, slide~29: in round~1 every $s'_i \ne s_i^D$ is strictly
  dominated by $s_i^D$) and hence the unique Nash equilibrium (Lecture~2, slide~17). Without
  strictness uniqueness can fail: if dominance only required ``$\ge$ against every $s_{-i}$'', a
  player whose payoffs are all equal would have every strategy dominant, and every profile of such
  players would be a dominant strategy equilibrium.
\end{itemize}

\interp A strictly dominant strategy is a best reply that does not depend on beliefs. Two such
strategies for the same player would each have to do strictly better than the other against the
same opponents, which is impossible. That is why the dominant-strategy prediction, when it exists,
is a single profile, and why it needs only rationality, not common knowledge of it.

\verify \texttt{check\_A1.py} draws 3{,}000 random finite games (two or three players, two to four
strategies each, integer payoffs), plants a strictly dominant strategy for every player in half of
them, and checks in each game that, whenever a strict dominant strategy equilibrium exists, no
other profile consists of weakly or strictly dominant strategies, and that it is the only IESDS
survivor. Tadelis gives no answer.

\mistakes
\begin{itemize}
  \item Proving only that there is no second \emph{strict} dominant strategy equilibrium. The
  proposition says ``dominant'', and the same two-line argument rules out weakly dominant ones.
  \item Comparing the two profiles $s^D$ and $s'$ as a whole. Dominance is a property of one
  player's strategy against all of the others' profiles; the contradiction comes from a single
  player whose strategies differ.
  \item Forgetting that $S_{-i}$ must be non-empty for the contradiction (one $\hat s_{-i}$ is
  needed). It always is in a game, but the step should be stated.
\end{itemize}

\answer{If $s^D$ is a strict dominant strategy equilibrium and $s' \ne s^D$, some player has
$s'_i \ne s_i^D$. Strict dominance of $s_i^D$ gives $v_i(s_i^D, s_{-i}) > v_i(s'_i, s_{-i})$ for
every $s_{-i}$, so $s'_i$ cannot be dominant, strictly or weakly (that would need
$v_i(s'_i, s_{-i}) \ge v_i(s_i^D, s_{-i})$ for every $s_{-i}$). Hence $s'$ is not a dominant
strategy equilibrium, and $s^D$ is the unique one.}
\end{exsol}
